\documentclass[a4paper,11pt]{article}

\usepackage[utf8]{inputenc}
\usepackage[portuguese]{babel}
\usepackage[T1]{fontenc}
\usepackage{geometry}
\usepackage{xcolor}
\usepackage{hyperref}
\usepackage{graphicx}

\geometry{margin=2.5cm}

\begin{document}

\title{Relatório de Estágio : Avaliação de Plataformas de Gestão de Dispositivos ESP32}
\author{Francisco José Nogueira da Rocha Matos}
\date{\today}

\maketitle

\tableofcontents

\section{Introdução}
Durante este estágio de 1 mês na Watgrid, avaliei plataformas e sistemas operativos para o ESP32, com o objetivo de identificar quais melhor suportam o caso de uso real da empresa. O foco inicial foi a atualização remota de firmware (OTA), alargado depois à facilidade de integração de \textit{drivers} de sensores I2C e comunicação LoRa. Foram avaliados o Zephyr RTOS e o ESP-IDF nativo (camada de firmware), o ESP RainMaker, o ThingsBoard e o Azure IoT Hub + DPS (plataformas de gestão de dispositivos na \textit{cloud}), e ainda o Mongoose OS como candidato adicional.

\section{Zephyr RTOS}
\begin{itemize}
    \item Conectividade Wi-Fi validada, com ligação automática programática.
    \item OTA via MCUboot+mcumgr (série): funcional, modo \textit{upgrade-only} (sem \textit{rollback}).
    \item OTA via LwM2M/Leshan (Wi-Fi): validado com hardware real, sem falhas -- 680KB, \textit{rollback} do MCUboot aplicado com sucesso. Requer um \textit{patch} adicional ao checkout (\textit{issue} \#77952/PR \#78121), ainda não incluído numa \textit{release} oficial.
    \item \textit{Drivers} de sensores I2C e LoRa: subsistema nativo \textit{in-tree} (\texttt{drivers/sensor/}, \texttt{drivers/lora/}). IIS2ICLX com \textit{driver} nativo (I2C+SPI); KX122 e MCP3221 não encontrados; LoRa nativo e completo (SX126x/SX127x/RYLR).
\end{itemize}

\section{ESP-IDF Nativo}
\begin{itemize}
    \item OTA via HTTP \textbf{sem qualquer falha}, com Wi-Fi ativo durante a escrita -- exatamente o cenário que falhava no Zephyr, confirmando que a causa é do \textit{port} Zephyr, não do hardware/rede.
    \item \textit{Rollback} automático confirmado: imagem defeituosa (\texttt{abort()}) reverte sozinha para a versão estável, com proteção contra repetição automática.
    \item Esforço de integração baixo, sem as dificuldades de Kconfig/\textit{devicetree} do Zephyr.
    \item \textit{Drivers} de sensores I2C e LoRa: sem subsistema nativo (só periféricos genéricos -- I2C/SPI/ADC interno). IIS2ICLX, KX122 e MCP3221 ausentes do ESP Component Registry; LoRa sem suporte nativo, mas com componentes de terceiros maduros (\texttt{radiolib}, \texttt{sx127x}).
\end{itemize}
\newpage
\section{ThingsBoard}
\begin{itemize}
    \item Ligação MQTT com \textit{access token}; telemetria, \textit{shared attributes} (config remota) e RPC (\texttt{getStatus}/\texttt{reboot}) validados.
    \item OTA próprio sobre MQTT (blocos com \textit{checksum} incremental) -- sem falhas.
    \item \textit{Rollback} confirmado; \textbf{achado}: proteção anti-repetição não vinha incluída, causando ciclo infinito download/falha/reversão -- corrigido com memória de falha em NVS.
    \item Provisionamento automático (credencial genérica $\rightarrow$ token individual) e alarme de desconexão com email, ambos validados.
    \item Lacuna de segurança identificada e corrigida: transporte migrado de \texttt{mqtt://} para \texttt{mqtts://}.
    \item Esforço de integração baixo, sem cartão de crédito/débito, grátis até 5 dispositivos.
    \item \textbf{PoC Zephyr + ThingsBoard}: ao contrário do RainMaker (SDK preso ao ESP-IDF), o ThingsBoard só precisa de um cliente MQTT -- escrito e validado com hardware real sobre o MQTT nativo do Zephyr (provisionamento, MQTTS, telemetria). Dois bugs reais corrigidos: config. \textit{default} do mbedTLS no Zephyr não tinha curvas elípticas (\texttt{MBEDTLS\_ECP\_C}) nem SHA-384 ativos, ambos exigidos pela cadeia de certificados do \textit{broker}.
\end{itemize}

\section{ESP RainMaker}
\begin{itemize}
    \item Esforço de integração inicial \textbf{mais baixo de todas as plataformas} (\textit{zero-config}, gestor de dependências automático).
    \item \textbf{Nota -- \textit{vendor lock-in}:} SDK proprietário, só funciona em ESP32/Espressif -- ao contrário do ThingsBoard/Azure (qualquer fabricante). Esforço mais baixo tem esta contrapartida.
    \item Provisionamento via BLE + app móvel obrigatória.
    \item OTA e \textit{rollback} ativos por omissão, sem código adicional.
    \item RPC customizado validado; \textit{Node Groups} (gestão de frota) confirmados só a nível conceptual (um único dispositivo).
    \item Achado espontâneo: notificação de desconexão \textit{push}, sem qualquer configuração prévia.
    \item \textbf{Reavaliação:} \texttt{Reboot} e \texttt{Wi-Fi Reset} nativos (\textit{System Service}, zero código, testados com sucesso via CLI/consola/app); RSSI e saúde de memória nativos via \textit{Insights} (\textit{dashboard} incluído, sem código nosso); versão de firmware exposta como metadado nativo. Resultado: \textbf{única das três plataformas com as 7 capacidades testadas nativas} (ThingsBoard e Azure só têm 1 cada).
    \item \textbf{Achado -- fragilidade do provisionamento:} reprovisionar a mesma placa (reutilizada de outros testes) falhou repetidamente via BLE/app oficial (\textit{ESPSessionError}), só resolvido com limpeza completa da flash -- e voltou a falhar numa tentativa seguinte.
\end{itemize}

\newpage
\section{Azure IoT Hub + Device Provisioning Service (DPS)}
Avaliação prática subsequente (exige cartão de crédito/débito).
\begin{itemize}
    \item \textbf{Provisionamento:} DPS + IoT Hub (\textit{tier} Free) configurados; SAS tokens calculados localmente (HMAC-SHA256), exigindo sincronização SNTP. Testado com sucesso tanto \textit{Individual} como \textbf{\textit{Group Enrollment}} (chave partilhada, cada dispositivo deriva a sua própria chave -- duas identidades distintas coexistentes confirmadas).
    \item \textbf{Direct Methods} (RPC) e \textbf{Device Twin} (config remota \textit{desired}/\textit{reported}) validados com sucesso.
    \item \textbf{OTA:} serviço gerido oficial indisponível no \textit{tier} Free -- sinalizado via Device Twin, reaproveitando o mecanismo HTTP do ESP-IDF nativo. Sucesso confirmado; teste de falha revelou e permitiu corrigir um \textbf{bug real} na proteção anti-ciclo (ordem de reconciliação pós-\textit{reboot}), com paralelo direto ao mesmo problema já encontrado no ThingsBoard.
    \item \textbf{Gestão de frota:} \textit{Configurations + Deployments} validado com dois dispositivos (um real, um simulado), superando o nível apenas conceptual do RainMaker.
    \item \textbf{Alarmes:} configuração explorada até ao fim, mas não ativada -- único recurso do PoC com custo mensal explícito (\$0.10 USD), decisão deliberada de âmbito.
    \item \textbf{C2D messages} (mensagem pontual assíncrona, cloud$\rightarrow$dispositivo) e \textbf{Message Routing} (filtragem de mensagens por condição) validados -- exploração lateral, mais sobre dados do que sobre gestão do dispositivo.
    \item Esforço de integração no firmware alto (protocolo do DPS e SAS tokens implementados à mão, sem SDK oficial pesado).
\end{itemize}

\section{Síntese Comparativa}

\subsection{OTA (dispositivo)}
\begin{center}
\resizebox{\textwidth}{!}{%
\begin{tabular}{|l|l|l|l|l|l|}
\hline
\textbf{Critério} & \textbf{MCUboot} & \textbf{LwM2M} & \textbf{ESP-IDF} & \textbf{ThingsBoard} & \textbf{RainMaker} \\
\hline
Transporte & Série & Wi-Fi/CoAP & Wi-Fi/HTTP & Wi-Fi/MQTT & Wi-Fi/HTTPS \\
\hline
Concluiu & Sim & Não & Sim & Sim & Sim \\
\hline
\textit{Rollback} & Não & Não testado & Confirmado & Confirmado & Por omissão \\
\hline
Esforço & Alto & Alto & Baixo & Baixo & Muito baixo \\
\hline
\end{tabular}%
}
\end{center}

\subsection{Plataformas de Gestão}
\begin{center}
\resizebox{\textwidth}{!}{%
\begin{tabular}{|l|l|l|l|l|}
\hline
\textbf{Critério} & \textbf{ESP-IDF} & \textbf{ThingsBoard} & \textbf{RainMaker} & \textbf{Azure DPS+Hub} \\
\hline
Hardware suportado & Só ESP32 & Qualquer & \textbf{Só ESP32/Espressif} & Qualquer \\
\hline
Dashboards & Não & Sim & Sim & Não testado \\
\hline
Config. remota & Não & Sim & Sim & Sim, confirmado \\
\hline
RPC & Não & Sim & Sim & Sim, confirmado \\
\hline
OTA via cloud & Não & Sim & Sim & Parcial  \\
\hline
\textit{Rollback} & Confirmado & Confirmado & Confirmado & Confirmado (bug corrigido) \\
\hline
Alarmes & Não & Sim & Sim (espontâneo) & Explorado, não ativado (custo) \\
\hline
Provisionamento & Não & Sim & Sim (app/BLE) & Sim (individual + grupo) \\
\hline
Gestão de frota & Não & Sim & Só conceptual & Sim, confirmado \\
\hline
Esforço inicial & Baixo & Baixo & Muito baixo & Alto \\
\hline
Custo & Grátis & Grátis (5 disp.) & Grátis (20 disp.) & Grátis (\textit{tier} Free) \\
\hline
\end{tabular}%
}
\end{center}

\subsection{Ver vs. Gerir: Quem Implementa}
Distinção relevante: separar \textbf{ver} (leitura de estado) de \textbf{agir} (comando), e se cada um vem pronto da \textbf{plataforma} ou foi feito \textbf{manualmente}.
\begin{center}
\resizebox{\textwidth}{!}{%
\begin{tabular}{|l|l|l|l|}
\hline
\textbf{Capacidade} & \textbf{ThingsBoard} & \textbf{RainMaker} & \textbf{Azure (DPS+Hub)} \\
\hline
Ver versão de firmware & Manualmente (telemetria) & \textbf{Plataforma} (metadado nativo do nó) & Manualmente (\textit{reported}) \\
\hline
Disparar atualização (OTA) & Híbrido (pacote: plataforma; download: manual) & \textbf{Plataforma} (consola própria, zero código) & Manualmente (Device Twin) \\
\hline
Reiniciar remotamente & Manualmente (RPC customizado) & \textbf{Plataforma} (\textit{System Service} nativo) & Manualmente  \\
\hline
Ver força de sinal Wi-Fi (RSSI) & Não testado & \textbf{Plataforma} (Insights, zero código) & Manualmente (\texttt{wifiRssi} no \textit{twin}) \\
\hline
Ver estado de conectividade & \textbf{Plataforma} (\texttt{active}, nativo) & \textbf{Plataforma} (notificação espontânea) & \textbf{Plataforma} (\texttt{connectionState}) \\
\hline
Agir sobre conectividade & Não testado & \textbf{Plataforma} (\textit{Wi-Fi Reset} nativo) & Não testado \\
\hline
Ver memória/saúde geral & Manualmente (RPC customizado) & \textbf{Plataforma} (Insights, zero código) & Manualmente (\textit{twin} + Direct Method) \\
\hline
\end{tabular}%
}
\end{center}
RainMaker: único com todas as 7 capacidades nativas (0 código nosso); ThingsBoard e Azure têm só 1 cada (conectividade).

\subsection{Drivers I2C e LoRa: Zephyr vs. ESP-IDF}
\begin{center}
\resizebox{\textwidth}{!}{%
\begin{tabular}{|l|l|l|}
\hline
\textbf{Critério} & \textbf{Zephyr} & \textbf{ESP-IDF} \\
\hline
Subsistema nativo de \textit{drivers} de sensores & Sim & Não \\
\hline
IIS2ICLX & \textbf{Nativo} & Não encontrado \\
\hline
KX122 / MCP3221 & Não encontrados & Não encontrados \\
\hline
LoRa nativo & \textbf{Sim} (SX126x/SX127x/RYLR) & Não (terceiros: radiolib, sx127x) \\
\hline
\end{tabular}%
}
\end{center}
Nenhum \textit{framework} tem os 3 sensores prontos, mas o Zephyr trata sensores/LoRa como parte do próprio SO (API unificada); no ESP-IDF depende sempre de bibliotecas de terceiros.

\textbf{Nota sobre módulos LoRa concretos (RN2483, WLR089U0):} são módulos com comando por texto via UART, não transcetores SPI \textit{nus} -- nenhum \textit{framework} tem isto pronto (RadioLib confirmadamente também não cobre o RN2483), mas é uma tarefa simples (\textit{wrapper} de UART, não driver de rádio), com bibliotecas de referência publicadas para adaptar.

\subsection{Prova de conceito: driver RN2483 (LoRa)}
Escrito e compilado com sucesso um segundo driver \textit{out-of-tree} para Zephyr, agora para o módulo RN2483 (comandos de texto via UART, não SPI). Liga-se à mesma API nativa de LoRa (\texttt{lora\_config}/\texttt{lora\_send}) -- do ponto de vista de quem usa, é intercambiável com um transcetor nativo. 302 linhas totais (mais que o MCP3221, por ter mais comandos de configuração de rádio), compilou à primeira, sem \textit{warnings}. Sem módulo físico, validado só ao nível de compilação/API, não de comunicação real.



\section{Mongoose OS}
Testado com o mesmo rigor prático (instalação, compilação, gravação real):
\begin{itemize}
    \item \textbf{Manutenção}: último \textit{release} dez/2021, último código real fev/2023 (correção ESP32-C3) -- commits recentes são só links no README.
    \item \textbf{O que o substituiu}: a Cesanta não desapareceu -- passou a concentrar-se na \textit{Mongoose Library} (\texttt{cesanta/mongoose}), uma biblioteca de rede em C (TCP/IP, HTTP, MQTT, TLS, agora com \textit{rollback} de OTA), usada \textbf{dentro} de outro sistema (bare-metal/FreeRTOS/Zephyr/ESP-IDF), não um SO próprio. Muito ativa: \textit{releases} a cada poucos meses e código novo até à véspera desta consulta. Âmbito bem mais estreito que o Mongoose OS (só rede, sem drivers/OTA de app/scripting) -- não avaliada em profundidade, fica como referência.
    \item \textbf{Teste prático}: instalação, compilação (\texttt{demo-c}) e gravação bem-sucedidas no hardware real -- \textbf{ainda funciona hoje}. Achado: \texttt{mos build} compila na nuvem da Cesanta, não localmente (dependência extra, ponto único de falha, ao contrário de Zephyr/ESP-IDF).
    \item \textbf{Drivers}: mesmo padrão de Zephyr/ESP-IDF -- KX122, IIS2ICLX e MCP3221 não têm biblioteca pronta; sem biblioteca de LoRa.
    \item \textbf{Features notáveis (não adotado, mas de registar)}: \textit{scripting} JavaScript sem recompilar (mJS), dashboard/OTA próprio (mDash), config remota nativa, RPC extensível (\texttt{RPC.List}), provisionamento pronto para AWS/Azure/GCP/Watson, hardware suportado (ESP32, ESP8266, STM32, CC3200/3220) -- nenhum framework avaliado tem equivalente.
\end{itemize}
\textbf{Conclusão}: não está avariado, só sem evolução e dependente de infraestrutura de nuvem de terceiros -- risco de manutenção, não limitação técnica.

\section{Conclusão}



O \textbf{RainMaker} é quem melhor cumpre o critério -- único com as 7 capacidades de gestão do dispositivo nativas, zero código nosso. \textbf{Mas} tem duas reservas: (1) o provisionamento via BLE/app mostrou-se pouco fiável na prática (falhas de sessão repetidas, só resolvidas com limpeza total da flash) -- risco operacional real para uma implantação com vários sensores em campo, sem garantia de técnico com telemóvel disponível sempre que preciso; (2) é um SDK proprietário, só ESP32/Espressif -- compromete a Watgrid a este fornecedor de hardware, sem a flexibilidade de diversificar chips que o ThingsBoard/Azure permitem.

O \textbf{ThingsBoard} continua a opção mais equilibrada se a fiabilidade/simplicidade do provisionamento (sem app/BLE) pesar mais do que a natividade capacidade-a-capacidade -- ao custo de implementar manualmente quase tudo. Reserva: limite de 5 dispositivos grátis.

A \textbf{Azure} (testada em profundidade nesta ronda) iguala ou supera o ThingsBoard tecnicamente -- incluindo \textit{Group Enrollment} e gestão de frota, que superam mesmo o RainMaker -- mas com o maior esforço de integração de firmware das três. Melhor candidato para escalar além do limite do ThingsBoard.

\textbf{Recomendação (gestão do dispositivo):} RainMaker como foco imediato (melhor no critério que importa agora), \textbf{condicionado} a validar a fiabilidade do provisionamento em escala real antes de comprometer. ESP-IDF + ThingsBoard como alternativa segura e já validada de ponta a ponta caso essa validação falhe.

\textbf{Firmware:} com LoRa confirmado como requisito crítico do projeto, e sendo nativo só no Zephyr (não no ESP-IDF), o Zephyr é candidato firme para a camada de firmware -- condicionado a aplicar um \textit{patch} adicional para OTA fiável sobre Wi-Fi, validado com hardware real, e a assumir o custo de \textit{drivers} próprios para KX122/MCP3221.

\textbf{Decisão a resolver:} o RainMaker corre sobre ESP-IDF, não sobre Zephyr -- se o LoRa nativo do Zephyr for requisito duro, isso é incompatível com o RainMaker. \textbf{Zephyr + ThingsBoard} deixou de ser só teórico: validado com hardware real (provisionamento, MQTTS, telemetria, sem SDK de firmware específico exigido) -- é a combinação recomendada nesse cenário.

\end{document}